%!TEX root = ../csuthesis_main.tex
\keywordsen{Hazardous Driving Scenarios \ \ Reinforcement Learning \ \ Scenario Generation \ \ CARLA Simulation}
\begin{abstracten}

With the rapid development of autonomous driving technology, how to effectively evaluate the safety of autonomous driving systems in complex traffic environments has become an important research issue. Traditional road testing is gradually being replaced by simulation testing due to its high cost, significant safety risks, and difficulty in reproducing special scenarios. To meet the requirements of autonomous driving safety testing, this paper proposes a classification and generation method for hazardous driving scenarios based on the CARLA simulation platform.
This paper combines manual annotation and reinforcement learning techniques to identify and classify different types of hazardous scenarios, and converts the generated scenarios into .xosc files that conform to the OpenSCENARIO standard. In addition, typical hazardous factors, including adverse weather conditions, traffic violations, and interactions among multiple traffic participants, are analyzed. On this basis, a semi-automatic data generation method combining manual annotation and Proximal Policy Optimization (PPO) reinforcement learning is adopted to construct a dataset containing ten types of high-risk driving scenarios, such as vehicle cut-in, lane changing, emergency braking, and pedestrian crossing.
To evaluate the generated scenarios, an evaluation framework is established from multiple aspects, including scenario reproducibility, hazard level, and repeatability, and the experimental results are analyzed through charts and visualizations. The results demonstrate that the proposed method can effectively improve the efficiency of hazardous scenario generation and produce realistic and complex traffic scenarios, thereby providing support for the safety testing of autonomous driving systems.

\end{abstracten}